\section{Related Work}\label{sec:related_work}

We survey prior research along four themes---short-horizon load forecasting, graph-based spatio-temporal modelling, multi-task learning in energy systems, and explainable artificial intelligence---drawing on a structured review of fifty-five studies (2023--2026) from the Bangladesh smart-grid programme. The focus is on methodologies, assumptions, and unresolved challenges that bound the Correlation-Aware Multi-Task Forecasting Framework (PF-STGT), not on enumerating individual publications.

\subsection{Load Forecasting}\label{subsection:load_forecasting}

Short-horizon electricity load forecasting remains the largest cluster in the reviewed literature. Methods range from classical regressors to deep recurrent, convolutional, and transformer architectures applied mainly to scalar or multivariate load series. Much of this work targets residential or AMI consumption at sub-daily resolution, assuming dense sampling rather than administratively defined regional nodes on a national footprint. Ensemble trees and gradient boosting remain strong tabular baselines when spatial structure is implicit in feature engineering \cite{Breiman2001RandomForest,Chen2016XGBoost,Musbah2025LoadForecasting}.

Deep learning studies emphasise temporal pattern extraction through stacked recurrent encoders, Autoformer- and transformer-style models \cite{Hochreiter1997LSTM,Wu2021Autoformer,Helmy2024AutoformerEV,Guillen2026LoadComposition}, and multi-encoder designs that fuse exogenous embeddings with load history. Common assumptions include stationary seasonality, long uninterrupted series, and horizons aligned to market cadences that may differ from daily national operational reporting. Developing-economy national grids with division-level evening-peak reporting and co-visible limitation indicators are effectively absent from the collected corpus.

A parallel stream addresses load shedding, under-frequency response, and emergency control rather than ex-ante day-ahead forecasting \cite{Paphitis2025SpatialLoadShedding,Wang2026SHAPLoadShedding}. These studies optimise shedding sequences, preserve stability, or manage islanded microgrids using simulation-based or event-centric datasets. Forecasting graded daily stress for planning is related to shedding research but distinct from optimising control actions at disturbance timescales.

Across this cluster, external-validity reporting is often weak---South Asian division-level case studies are rare, conference abstracts can be metadata-sparse, and evaluation metrics are sometimes mixed without clear documentation. Open challenges include adapting AMI-centric pipelines to daily national timelines, modelling regional heterogeneity where one division dominates load, and coupling load forecasts with system-pressure indicators derived from limitation dynamics.

\subsection{Graph-based Spatio-Temporal Forecasting}\label{subsection:graph_based_spatio_temporal_forecasting}

Graph neural networks and graph-attention architectures \cite{Kipf2017GCN,Sun2024GLFN} construct nodes from geographic areas, charging stations, or feeders, with edges defined by administrative proximity, infrastructure topology, or correlation in historical series. Spatio-temporal graph convolutional networks propagate information across fixed neighbourhoods before or after temporal encoding \cite{Shi2024EVChargingGNN,Zhuang2024MultiScaleGNN,Wu2023TransferLearningGNN}.

Transformer-based spatio-temporal forecasting stacks graph attention with temporal transformers or applies dual-channel spatial--temporal decoders \cite{Zhao2024GraphAttentionTransformer,Bentsen2023GraphTransformer,Vaswani2017Attention}, typically on US or European residential benchmarks with spatial semantics unlike a nine-division national grid. Geographic adjacency is a common inductive bias; data-driven correlation graphs appear less often as the sole spatial prior, and hybrid geographic--correlation constructions are seldom ablated against correlation-only alternatives on developing-economy records.

Dataset assumptions diverge from the Bangladesh setting: node counts, sub-hourly resolution, and rare co-forecasting of stress with regional demand in the same graph-transformer pipeline. Shallow graph-convolution baselines \cite{Yu2018STGCN,Zhao2020TGCN} test whether graph propagation alone suffices when temporal encoding is limited. Reporting of graph-construction splits is often insufficient, explainability of learned spatial weights remains sparse, and evidence is limited on when geographical priors help versus add noise on strongly correlated national demand series. Selecting graph priors when inter-regional coupling reflects shared national growth rather than border adjacency, determining whether transformer capacity is conditional on graph topology, and extending graph-temporal models to division-level nodes with synchronised limitation covariates remain open problems. None of the high-relevance collected studies address Bangladesh shedding context or daily OSI-style stress targets.

\subsection{Multi-Task Learning}\label{subsection:multi_task_learning}

Multi-task learning in energy systems \cite{Caruana1997Multitask} occupies a narrow band relative to single-task load forecasting and shedding control. Reported formulations couple related energy vectors or hierarchical horizons within a single sector \cite{Song2024MultiTaskEnergy,Zhuang2024MultiScaleGNN}. Explicit coupling of continuous multi-node regional demand with a daily operational stress index derived from shedding intensity, reserve margin, and limitation pressure is not represented among collected comparators.

Task isolation remains dominant: demand papers treat load as the sole target; shedding papers optimise control; operational-stress studies focus on transmission assets rather than forecastable daily stress scalars. Auxiliary targets, where present, are not paired with graph-level spatial encoders on national regional footprints. The literature offers limited guidance on balancing megawatt-scale robust regression against bounded stress targets or on repairing stress-head collapse when naive multi-task weighting yields uninformative auxiliary outputs.

Scarce ablation of single-task versus multi-task ceilings, infrequent reporting of Pareto trade-offs when auxiliary tasks marginally degrade primary metrics, and absence of developing-economy national case studies with documented chronological splits further limit transferability. Open questions include when multi-task learning is preferable to separate pipelines, how to calibrate task weights under composite-target semantics, and whether shared graph-temporal encoders transfer stress-relevant structure beyond single-task models with global limitation covariates.

\subsection{Explainable AI}\label{subsection:explainable_ai}

Explainability for load and smart-grid forecasting is sparsely represented \cite{Baur2024XAIReview}. Where it appears, it clusters into SHAP attributions for control or shedding policies \cite{Wang2026SHAPLoadShedding,Lundberg2017SHAP}, interpretable ML frameworks for decentralised prediction \cite{Cifci2025InterpretableSmartGrid}, and review articles cataloguing feature-importance practices for scalar load series without graph-temporal multi-task outputs.

Attribution is often treated as a supplement to control validation rather than operator-facing diagnostics for coupled regional demand and national stress forecasts. Graph-temporal models rarely export spatial attention, coalition-grouped gradients, and dual-path comparison against engineered stress components in one evaluation programme. Review-level work \cite{Baur2024XAIReview} acknowledges that explainability methods can disagree on feature rankings, yet few studies document that disagreement for multi-node graph transformers.

Attributions are frequently computed on small validation samples without operator studies, post-hoc sensitivity is sometimes conflated with physical causation, and stress and demand explanations are seldom jointly audited. Standardising multi-method attribution for dual-task dashboards, aligning explanations with engineered stress components, and extending microgrid-oriented SHAP workflows to national graph models with eleven feature coalitions remain unresolved.

\subsection{Summary and Research Gap}\label{subsection:summary_and_research_gap}

The reviewed literature offers capable methods within each thematic silo but seldom integrates them for the problem addressed here. Four cross-cutting gaps emerge.

First, geographic validity is limited: no collected peer-reviewed spatio-temporal study uses Bangladesh division-level smart-grid daily data \cite{Hossain2025BangladeshSmartGrid}, and developing-economy national grids remain under-represented. Second, task formulation remains fragmented: continuous regional demand is rarely coupled with operational-stress targets in one multi-task graph-temporal model, while shedding research predominantly targets control rather than day-ahead planning visibility. Third, methodological integration is incomplete: only a minority combine graph-based spatial coupling with transformer-style temporal modelling, and few co-design limitation-aware covariates with spatio-temporal architectures or document chronological evaluation with ablations and corrected inferential testing. Fourth, explainability is rarely embedded in coupled demand-and-stress graph forecasts with explicit multi-method reporting.

These gaps motivate PF-STGT as a targeted response rather than a claim of universal superiority. The framework predicts joint next-day regional evening-peak demand and graph-level OSI on Bangladesh operational records under leakage-safe chronological evaluation, using a train-estimated correlation graph as an empirically testable spatial prior, parallel-fusion graph-aware and temporal transformer encoding, and a repaired multi-task loss that stabilises stress learning. Post-hoc attribution---coalition-level gradients, permutation comparisons, attention export, and case-stratified dual-path stress review---is treated as a reporting obligation. The present study sits at the intersection of Bangladesh division-level analytics, multi-task demand-and-stress coupling, graph-temporal modelling with documented graph-prior ablation, and explainability on frozen multi-task outputs. Methodology and empirical evaluation follow in subsequent sections.
